% !PRINTED 442
\weberpage{442}
%<<CONT

\begin{tabular}{@{}l l l l@{}}
1. & Substitutionen mit drei Polen, & & \\
2. & \textquotedbl & \textquotedbl & zwei Polen,\\
3. & \textquotedbl & \textquotedbl & einem Pol,\\
4. & \textquotedbl & \textquotedbl & einer Hauptaxe und einem Pol,\\
5. & \textquotedbl & \textquotedbl & einer Hauptaxe.\\
\end{tabular}

Als letzter Fall w\"urde noch die Multiplication aufzuz\"ahlen sein, f\"ur den jeder beliebige Punkt sein eigenes Bild ist.

Die Art der Substitution bleibt erhalten bei jeder Transformation.

Wenn ein Punkt unge\"andert bleibt durch eine Substitution $A$, so bleibt er auch bei jeder Wiederholung von $A$, also bei $A^2, A^3,\dots$, unge\"andert. Die Pole von $A$ finden sich daher immer unter den Punkten, die bei der Abbildung durch $A^2$ ihre eigenen Bilder sind. Es kann aber wohl der Fall eintreten, dass z.\,B. $A$ drei Pole hat, w\"ahrend eine Potenz, $A^k$, eine Hauptaxe besitzt. Dann m\"ussen zwei der Pole von $A$ auf der Hauptaxe von $A^k$ liegen, und der dritte Pol von $A$ ist der einzelne Pol von $A^k$. Ist n\"amlich $A$ eine Substitution mit drei Polen, so k\"onnen wir sie, indem wir die Pole in die Ecken des Coordinatendreiecks legen, in die Normalform
\[
A=\begin{pmatrix}\alpha, & 0, & 0\\ 0, & \beta, & 0\\ 0, & 0, & \gamma\end{pmatrix}
\]
setzen, worin $\alpha,\beta,\gamma$ von einander verschieden sind. Ist nun $\beta^k=\gamma^k$, also $\beta$ von $\gamma$ um eine $k^{\text{te}}$ Einheitswurzel als Factor unterschieden, so hat $A^k$ die Linie $x_1=0$ zur Axe, und die gegen\"uberliegende Ecke des Coordinatendreiecks zum Pol.

Ist $A$ von endlichem Grade $\mu$, so sind $\alpha,\beta,\gamma$ $\mu^{\text{te}}$ Einheitswurzeln. Ist $A^k$ die niedrigste Potenz von $A$, die eine Hauptaxe hat, so muss $\beta:\gamma$ primitive $k^{\text{te}}$ und zugleich $\mu^{\text{te}}$ Einheitswurzel sein, und folglich muss $k$ ein echter Theiler von $\mu$ sein; wenn $k$ relativ prim zu $\mu$ ist, so hat $A^k$ dieselben drei Pole wie $A$.

\begin{center}
\S.\,117.

\gesperrt{Anwendung auf die Gruppe $G_{168}$. Siebenz\"ahlige Pole.}
\end{center}

Wir wollen nun die Pole und Axen der Substitutionen unserer Gruppe $G_{168}$ aufsuchen.
